\documentclass{article}
\usepackage[left=2cm, right=2cm, top=2cm, bottom=2cm]{geometry}
\usepackage{graphicx}
\usepackage{amsmath}
\usepackage{amssymb}
\usepackage{amsthm}
\usepackage{fancyhdr}
\usepackage{verbatim}
\usepackage{listings}
\usepackage{xcolor}
\usepackage{pdfpages}
\usepackage{cancel}
\usepackage{physics}
\usepackage{esint}
\usepackage{braket}

\lstset{
    language=C++,
    basicstyle=\ttfamily\footnotesize,
    keywordstyle=\color{blue},
    commentstyle=\color{green},
    stringstyle=\color{red},
    numbers=left,
    numberstyle=\tiny\color{gray},
    frame=single,
    breaklines=true,
    captionpos=b,
}


\title{PHYS 427: Lecture \# 1}
\author{Cliff Sun}

\newtheorem{theorem}{Theorem}[section]
\newtheorem{lemma}[theorem]{Lemma}
\newtheorem{definition}[theorem]{Definition}
\newtheorem{conjecture}[theorem]{Conjecture}
\newtheorem{proposition}[theorem]{Proposition}
\newtheorem{corollary}[theorem]{Corollary}
\newtheorem{one minute paper}[theorem]{One Minute Paper}

\pagestyle{fancy}
\lhead{\textbf{\thepage}\ \ \nouppercase{\rightmark}}
\chead{PHYS 427: Lecture \# 1}
\rhead{Cliff Sun}

\begin{document}

\maketitle

\section*{Lecture Span}
\begin{enumerate}
    \item Probability and Multiplicity
\end{enumerate}

\section*{Notes}
\begin{enumerate}
    \item Homework due Mondays at 11am in the homework box in interpass (427)
    \item Discussion is on Tuesdays, weekly quizzes
    \item 2 Midterms (October, November) in class, One final
\end{enumerate}

\section*{Microstates and Macrostates}

\begin{definition}
    A \underline{Microstate} is the complete description of every particle or configuration of a system, down to every detail. 
\end{definition}

\begin{definition}
    A \underline{Macrostate} is a description of a few measurable properties of a configuration. 
\end{definition}

In this class, we will review gas, solids, liquids, and magnets. 

\textbf{Example 1:} Consider a few gas molecules in a box, then each particle has a position and a momentum. Then, a microstate will be a description of all the atoms position and momentum, i.e. 

\begin{equation}
    \text{Microstate} = \{x_i, p_i\}
\end{equation}

The macrostate will be 

\begin{equation}
    \text{Macrostate} = p, T, V
\end{equation}

Or pressure, temperature, and volume. 


\textbf{Example 2:} Consider a polymer that is a long chain of several other molecules. Then 

\begin{equation}
    \text{Microstate} = \{\text{configuration of each polymer}\}
\end{equation}

\begin{equation}
    \text{Macrostate} = \{\text{Distance from end to end of polymer chain}\}
\end{equation}

\textbf{Example 3:} Consider a magnet with spins up and down in various configurations. This is an Ising Model. Then, the microstate would be 

\begin{center}
    Microstate $=$ Configuration of every spin in the magnet
\end{center}

\begin{center}
    Macrostate $=$ Average magnetization (average of all the spin values), $M$
\end{center}

\subsection*{Example}

The most probable macrostate is the one who has the most degenerate microstates. Consider a box divided into a left half and a right half. 
Consider the fact that we have $N$ particles in this box. Let the Macrostate be the number of particles in the left half $n_L$. Then $n_R = N - n_L$. 

\textbf{Microstate:} Configuration of $N$ particles as a list of $L/R$ state of each particle. That is, whether a particle is in the left or right half. Like 

\begin{center}
    LRRLRRL$\dots$
\end{center}

Consider multiplicity $\Omega(n_L, N)$. Then consider a table with the Microstate on the left column, $\Omega$ in the middle, and Macrostate on the right. 

Here, $N=1$
\begin{center}
\begin{tabular}{|c c c|} 
 \hline
 Microstate & $\Omega$ & Macrostate \\ [0.5ex] 
 \hline\hline
 Particle on the right & 1 & $n_L = 0$ \\ 
 \hline
Particle on the left & 1 & $n_L = 1$ \\
\hline
\end{tabular}
\end{center}
Total number of microstates is $2 = 2^1$, and total number of macrostates is $1$. 

For $N=2$
\begin{center}
\begin{tabular}{|c c c|} 
 \hline
 Microstate & $\Omega$ & Macrostate \\ [0.5ex] 
 \hline\hline
 2 particles on the right & 1 & $n_L = 0$ \\ 
 \hline
Symmetric particle distribution & 2 & $n_L = 1$ \\
\hline
2 particles on the left & 1 & $n_L = 2$ \\
\hline
\end{tabular}
\end{center}
Total number of microstates is $4 = 2^2$, and macrostates is $3 = 2 + 1$. 

For $N=3$,
\begin{center}
\begin{tabular}{|c c c|} 
 \hline
 Microstate & $\Omega$ & Macrostate \\ [0.5ex] 
 \hline\hline
 3 particles on the right & 1 & $n_L = 0$ \\ 
 \hline
1 particle on the left & 3 & $n_L = 1$ \\
\hline
1 particle on the right & 3 & $n_L = 2$ \\
\hline
3 particles on the left & 1 & $n_L = 3$ \\
\hline
\end{tabular}
\end{center}
Total number of microstates is $8 = 2^3$, and macrostates is $4 = 3 + 1$. 

\begin{definition}
    A \underline{Statistical Ensemble} is the total collection of all possible microstates of a system.
\end{definition}

The total number of microstates is $2^N$ and the total number of macrostates is $N+1$. 
This is for a system with only \underline{two states.}

Note, for small $N$, the probability of measuring a specific macrostate will not be that much different from each other. All the tools we develop in this course work in the \underline{Thermodynamic limit}, that is $N \rightarrow \infty$. 

Next, we derive the general formula for the multiplicity. For the first particle, we have $N$ choices. For the 2nd particle, we have $N-1$ choices. Etc. In general, for the $n-th$ L particle, we have $N- n_L + 1$ choices. 

Then, the number of choice is $N(N-1)(N-2)\dots(N - n_L + 1)$ if there are only $L$ slots.  

Here, the general formula for the multiplicity is 

\begin{equation}
    \Omega(n_L, N) = \frac{N(N-1)\dots(N - n_L +1)}{n_L!} = \frac{N!}{n_L!(N - n_L)!} = \begin{pmatrix}
        N \\
        n_L
    \end{pmatrix}
\end{equation}

This is also $N$ "choose" $n_L$. 

The reason why we divide the number of choices by $n_L$ is because, since all the particles are identical, we can rearrange them in $n_L!$ possible configurations in the left box, but we would get back the same configuration since all the particles are identical. 

Moreover, using the binomial thereom which states 

\begin{equation}
    (p + q)^N = \sum^N_{n=0}\begin{pmatrix}
        N \\
        n
    \end{pmatrix}p^{n}q^{N-n}
\end{equation}

\begin{equation}
    2^N = \sum \begin{pmatrix}
        N\\
        n_L
    \end{pmatrix} = \sum \Omega(n_L, N)
\end{equation}

\textbf{Fundamental Assumption:} In a closed, isolated system (i.e., non-interacting with surroundings), each microstate in a system is equally likely. Therefore, the probability of measuring one particle microstate is $1/2^N$. But some macrostates are more likely than others. 

Therefore, the probability that $n_L$ out of $N$ particles are to the left is 

\begin{equation}
    P(n_L, N) = \frac{\Omega(n_L, N)}{\sum_{n_L}\Omega(n_L, N) \iff \text{ total \# of states}}
\end{equation}

As $N \rightarrow \infty$, $P(n_L, N) \rightarrow \delta_{N/2}$, which is the delta function located at $N/2$ particles on the left and right. Check that 

\begin{align}
    \sum P &= 1\\
    \langle n_L \rangle &= \frac{N}{2} \\
    \langle f(n_L) \rangle &= \sum f(n_L)P(n_L)
\end{align}



\end{document}